\newpage
\section{Acta Número 16}
\label{sec:acta16}

\noindent \textbf{Fecha de la sesión:}\par
Lunes, 27 de Abril de 2026

\vspace*{0.5cm}
\noindent \textbf{Datos de la Sesión:}
\vspace{-0.2cm}
\begin{itemize}[noitemsep]
    \item \textbf{Modalidad:} Virtual - Google Meet.
    \item \textbf{Hora de Inicio:} 18:00
    \item \textbf{Hora de Fin:} 20:00
\end{itemize}

\vspace*{0.5cm}
\noindent \textbf{Orden del Día:}
\vspace{-0.2cm}
\begin{itemize}[noitemsep]
    \item Control de Asistencia.
    \item Sesión de \textit{Backlog Refinement} (Grooming) para el Sprint 3 (Epic 4).
    \item Priorización de Backlog e integración de deuda técnica (Sprints 1 y 2).
    \item Definición de plazos de entrega y acuerdos de cierre.
    \item Firmas y Aprobación de los Presentes.
\end{itemize}

\vspace*{0.5cm}
\noindent \textbf{Socios Fundadores Presentes:}
\vspace{-0.2cm}
\begin{enumerate}[noitemsep]
    \item Pardo Pozo Juan Diego (Jefe de Proyecto).
    \item Arnez Jaimes Mauricio.
    \item Quispe Flores Camila Belen.
    \item Ariñez Bautista Jim Gabriel.
    \item Medina Rodríguez Mateo Jhoustin.
    \item Molina Maldonado Tomas Manuel.
    \item Gutierrez Guzmán Angheline.
\end{enumerate}

\newpage
\subsection{Desarrollo del Acta}

\noindent \textbf{Control de Asistencia del Team}\par
Se procedió a la toma de asistencia a la hora acordada, corroborando la presencia de los 7 socios fundadores de la grupo-empresa Byte Busters S.R.L. Al contar con el quórum necesario y la disposición absoluta de todos los miembros, se dio inicio formal a la sesión de trabajo.

\vspace*{0.5cm}
\noindent \textbf{Sesión de \textit{Backlog Refinement} (Grooming) - Sprint 3 (Epic 4)}\par
Con la participación activa de todo el equipo, se llevó a cabo el proceso de \textit{Grooming} enfocado en las historias de usuario correspondientes a la Épica 4 (Sprint 3). Se realizó la descomposición detallada de dichas historias, logrando una clarificación profunda de los requerimientos de negocio y técnicos. En consenso, se establecieron de forma rigurosa los Criterios de Aceptación para garantizar el cumplimiento del \textit{Definition of Done} (DoD). Asimismo, el equipo en conjunto procedió con la estimación de esfuerzo de cada tarea, identificando y documentando las dependencias técnicas e interbloqueos para mitigar cualquier riesgo durante la fase de ejecución.

\vspace*{0.5cm}
\noindent \textbf{Priorización de Backlog e Integración de Deuda Técnica (Sprints 1 y 2)}\par
Como parte de la inspección y adaptación continua del producto, el equipo analizó el estado de los entregables anteriores. Se determinó, mediante acuerdo unánime, la necesidad técnica y funcional de corregir y finalizar historias de usuario pendientes provenientes del Sprint 1 y Sprint 2. Por consiguiente, se reordenó el \textit{Product Backlog}, priorizando estratégicamente estas tareas de corrección e integrándolas directamente al \textit{Sprint Backlog} del Sprint 3. Esta decisión busca asegurar la calidad y estabilidad de la arquitectura base antes de avanzar con los nuevos desarrollos.

\vspace*{0.5cm}
\noindent \textbf{Definición de Plazos de Entrega y Acuerdos de Cierre}\par
Para garantizar el flujo de valor y el cumplimiento de los objetivos del ciclo, el equipo en pleno acordó un límite de tiempo estricto para saldar la deuda técnica. Se estableció formalmente que todas las historias y tareas correspondientes a las correcciones del Sprint 1 y Sprint 2 deberán estar concluidas, validadas y cerradas a más tardar el día \textbf{sábado, 2 de mayo de 2026, a las 18:00 horas}. Todos los miembros confirmaron su entendimiento y compromiso total con este hito temporal.

\vspace*{0.5cm}
\noindent \textbf{Aprobación Oficial}\par
Habiendo estructurado exitosamente el Backlog para el tercer sprint, clarificado los requerimientos y establecido los plazos definitivos para las correcciones pendientes, todos los socios manifestaron su conformidad absoluta con las decisiones tomadas y la planificación estratégica trazada. Sin más puntos que tratar, se da por aprobada la presente acta y se procede con la firma correspondiente de todos los presentes.

\newpage
\noindent \textbf{Firmas y Aprobación de los Presentes}
\begin{center}
    \noindent
    \setlength{\tabcolsep}{0pt}
    \begin{tabular}{ccc}
        \espacioFirma{assets/img/firmas/mauricio.jpg}{Mauricio Jaimes Arnez}{13751221} & 
        \espacioFirma{assets/img/firmas/camila.jpg}{Camila Belen Quispe Flores}{13097244} &
        \espacioFirma{assets/img/firmas/jim.jpg}{Jim Gabriel Ariñez Bautista}{9517642} \\[1.0cm]
        \espacioFirma{assets/img/firmas/mateo.jpg}{Mateo Medina Rodríguez}{9453809} &
        \espacioFirma{assets/img/firmas/tomas.jpg}{Tomas Molina Maldonado}{8051701} &
        \espacioFirma{assets/img/firmas/angheline.jpg}{Angheline Gutierrez Guzmán}{13751815} \\[1.0cm]
    \end{tabular}
    \vspace{0.2cm} 
    \begin{tabular}{cc}
        \espacioFirma{assets/img/firmas/diego.jpg}{Juan Diego Pardo Pozo}{12747783}
    \end{tabular}
\end{center}

\vspace{0.5cm}
\noindent \textbf{Fecha de última revisión:} Lunes, 27 de Abril de 2026